\chapter{Sistemi lineari}

\section{Sistemi lineari}
Sia $\mathbf{x}=(x_1,\dots,x_n)^T\in \mathbb{R}^n$ un vettore colonna. Nel seguito verranno considerate le
seguenti norme di vettore:
\begin{itemize}
    \item $\|\mathbf{x}\|_1=\displaystyle\sum_{i=1}^n\left|x_i\right|$;
    \item $\|\mathbf{x}\|_2=\sqrt{\displaystyle\sum_{i=1}^nx_i^2}=\sqrt{\mathbf{x}^T\mathbf{x}}$,
    \quad \textbf{norma euclidea};
    \item $\|\mathbf{x}\|_\infty=\displaystyle\max_{1\leq i\leq n}\left|x_i\right|$.
\end{itemize}

Sia $\mathbf{A}=(a_{ij})_{i=1,\dots,m,\,j=1,\dots,n}\in\mathbb{R}^{m,n}$ una matrice di dimensioni $m\times n$.
Nella trattazione dei sistemi lineari verranno utilizzate le seguenti norme di matrice:
\begin{itemize}
    \item $\|\mathbf{A}\|_1=\displaystyle\max_{1\leq j\leq n}\sum_{i=1}^m\left|a_{ij}\right|$;
    \item $\|\mathbf{A}\|_2=\sqrt{\rho(\mathbf{A}^T\mathbf{A})}$\textbf{ norma spettrale}, ove
    $\rho(\mathbf{B})=\displaystyle\max_{1\leq j\leq n}\left|\lambda_i\right|$, con $\lambda_i$
    autovalore di $\mathbf{B}$, è detto \textbf{raggio spettrale} di $\mathbf{B}$;
    \item $\|\mathbf{A}\|_\infty=\displaystyle\max_{1\leq i\leq n}\sum_{j=1}^m\left|a_{ij}\right|$,\qquad
    $\left(\|\mathbf{A}\|_\infty=\|\mathbf{A}^T\|_1\right)$.
\end{itemize}

\begin{definition}{Norme compatibili}
    Data una norma di matrice e una di vettore, si dice che le due norme sono \textbf{compatibili} se
    \[\|\mathbf{Ax}\|\leq\|\mathbf{A}\|\,\|\mathbf{x}\|\]
    per ogni matrice $\mathbf{A}\in\mathbb{R}^{n,n}$ e un vettore $x\in\mathbb{R}^n$.
\end{definition}

\begin{observation}{Norme $1$, $2$ e $\infty$}
    Le norme precedentemente definite $1$, $2$ e $\infty$ di matrice e di vettore sono compatibili.
\end{observation}

\begin{observation}{Relazione tra raggio spettrale e norma}
    Per ogni norma di matrice compatibile con una norma di vettore, vale la seguente relazione
    \[\rho(\mathbf{A})\leq\|\mathbf{A}\|\]
\end{observation}

\begin{definition}{Matrice diagonale dominante per righe}
    Una matrice $\mathbf{A}\in\mathbb{R}^{n,n}$ si dice a \textbf{diagonale dominante per righe} se
    \[\left|a_{ii}\right|>\sum_{\substack{j=1\\j\ne i}}^n\left|a_{ij}\right|,\quad i=1,\dots,n\]
\end{definition}

\begin{definition}{Matrice diagonale dominante per colonne}
    Una matrice $\mathbf{A}\in\mathbb{R}^{n,n}$ si dice a \textbf{diagonale dominante per colonne} se
    \[\left|a_{jj}\right|>\sum_{\substack{i=1\\i\ne j}}^n\left|a_{ij}\right|,\quad j=1,\dots,n\]
\end{definition}

\begin{definition}{Matrice definita positiva}
    Una matrice \textbf{simmetrica $\mathbf{A}$} si dice \textbf{definita positiva} se
    \[\mathbf{x}^T\mathbf{Ax}>0,\quad\forall \mathbf{x}\ne\mathbf{0};\]
    oppure, equivalentemente, se
    \[\lambda_1(\mathbf{A})>0,\quad\forall i=1,2,\dots,n\]
    con $\lambda_i(\mathbf{A})$ autovalore di $\mathbf{A}$; oppure, equivalentemente, se
    \[\det(\mathbf{A}_k)>0,\quad\forall k=1,2,\dots,n\]
    ove $\mathbf{A}_k$ è la matrice di ordine $k$ formata dall'intersezione delle prime $k$ righe e
    $k$ colonne di $\mathbf{A}$.
\end{definition}

\begin{definition}{Matrice ortogonale}
    Una matrice $\mathbf{A}\in\mathbb{R}^{n,n}$ si dice \textbf{ortogonale} se le sue colonne (righe) formano
    un sistema ortonormale; in questo caso, denotata con $\mathbf{I}$ la matrice identità, si ha $\mathbf{A}^T
    \mathbf{A}=\mathbf{A}\mathbf{A}^T=\mathbf{I}$, quindi $\mathbf{A}^{-1}=\mathbf{A}^T$.
\end{definition}

\begin{definition}{Matrice di permutazione}
    Si definisce \textbf{matrice di permutazione}, una matrice ottenuta permutando le righe della matrice identità.
\end{definition}

\begin{observation}{Matrice di permutazione}
    Le matrici di permutazione in ogni riga e in ogni colonna hanno un solo elemento diverso da zero e uguale a $1$.\\
    Le matrici di permutazione sono matrici ortogonali.\\
    Come mostra il seguente esempio, se si moltiplica una matrice $\mathbf{A}$ o un vettore $\mathbf{b}$
    a \textbf{sinistra} per un'opportuna matrice di permutazione $\mathbf{P}$, si possono realizzare scambi
    di righe in $\mathbf{A}$ oppure id componenti in $\mathbf{b}$.
    \[\mathbf{PA} = \left(
        \begin{matrix}
            0 & 0 & 1 \\
            0 & 1 & 0 \\
            1 & 0 & 0
        \end{matrix}
    \right)\left(
        \begin{matrix}
            a_{11} & a_{12} & a_{13} \\
            a_{21} & a_{22} & a_{23} \\
            a_{31} & a_{32} & a_{33}
        \end{matrix}
    \right)=\left(
        \begin{matrix}
            a_{31} & a_{32} & a_{33} \\
            a_{21} & a_{22} & a_{23} \\
            a_{11} & a_{12} & a_{13}
        \end{matrix}
    \right)\]
\end{observation}
